\section{引言}
信息物理系统（Cyber-Physical Systems, CPS）~\cite{baheti2011cyber} 将计算能力与物理能力相结合，如今已广泛应用于自动驾驶和医疗设备等安全关键应用中~\cite{rajkumar2010cyber}。对于这类系统，任何故障都可能导致灾难性后果。线性时序逻辑（Linear Temporal Logic, LTL）是一种用于定义安全运行复杂规则的形式语言，例如要求及时响应或禁止不安全条件~\cite{pnueli1977temporal}。因此，验证 CPS 是否遵循其 LTL 规约的能力，对于在这些系统中建立信任并确保安全至关重要。

尽管意义重大，但对 CPS 的 LTL 规约进行形式化验证仍是一项重要挑战。主要困难来自于需要同时处理 LTL 性质的无限时域特征~\cite{baier2008principles}，以及物理系统的无限状态空间和复杂动力学。基于穷尽状态探索的传统模型检测技术，由于这类系统具有连续状态空间，并不能直接适用。基于证书的验证提供了一种有前景的替代方案。该方法不枚举状态，而是寻找一个满足某些不等式条件的数学函数，称为\emph{证书}，从而作为系统满足其规约的证明~\cite{prajna2004safety}。例如，屏障证书是一个函数，其水平集将可达状态与不安全状态分离，从而在不显式计算可达集的情况下证明安全性。证书设计中的一个核心关切是\emph{限制性}：某种特定证书类型施加的充分条件可能过于严格，以至于无法验证那些事实上满足规约的系统。另一个关切是\emph{合成效率}，因为证书参数必须自动寻找，并且计算成本会随搜索空间维度增加而增长。针对 LTL 验证，研究人员发展了基于证书的方法，包括状态三元组方法~\cite{wongpiromsarn2015automata} 和闭包证书~\cite{murali2024closure}。状态三元组方法为动力系统的 LTL 规约验证提供了一种自动机理论方法。该技术构造屏障证书，以阻止到达表示否定规约的自动机中的接受状态。由于它依赖单一证明策略，即通过屏障证书使乘积系统中的所有接受状态不可达，因此该方法可能具有较强限制性。因此，只要任何接受状态可从初始集到达，该方法在结构上就会失败，即使系统满足 LTL 规约，并且实际中接受转移只会触发有限多次。该方法的优势在于其与自动机结构的直接联系，但对于复杂规约，它可能需要大量证书。闭包证书通过转移不变式作用于系统状态对，从而提供了一种广义框架。该方法建立沿轨迹施加下降条件的证书，确保接受状态只被访问有限多次。该方法通过允许更一般的证书结构提供了更大的灵活性，并自然处理活性规约。虽然这些方法具有重要意义，但它们各自存在权衡。状态三元组方法可能具有限制性，而闭包证书由于必须探索较大的搜索空间而合成效率较低。这些局限性清楚表明，需要一种既能提供更大灵活性又不引入高计算成本的新框架。

为弥合这一差距，我们提出了一种通过证明证书验证 LTL 规约的新框架。我们的方法将 LTL 验证问题的重点转向否定式非确定性 Büchi 自动机（Nondeterministic Büchi Automaton, NBA）~\cite{hahn2020reward} 的转移。在我们的框架中，验证通过寻找两类互补证书来实现。转移安全性证书证明某些接受转移永不发生，而转移持久性证书保证其他接受转移只发生有限多次。通过将二者结合，我们能够形式化地保证不存在接受运行。为了寻找这些证书，我们基于反例引导归纳合成（Counterexample-Guided Inductive Synthesis, CEGIS）框架开发了一种自动化合成方法。该方法同时支持用于提升效率的多项式模板，以及用于增强表达能力的神经网络模板。本文的主要贡献总结如下：
\begin{itemize}

  \item 我们为具有连续状态空间的离散时间动力系统提出了一种新的 LTL 验证框架，该框架采用转移证书减少搜索空间，从而相较于闭包证书实现更高的计算效率。
  

  \item 我们引入了基于多项式模板和神经网络模板的转移证书自动化合成方法。这些模板在表达能力与计算效率之间取得平衡，使得能够针对从线性动力学到非线性动力学的多种系统类别高效生成证书。
  

  \item 我们通过案例研究展示了所提出证书的有效性，包括 Kuramoto 振子和双房间温度模型。实验结果表明，转移证书实现了显著加速，突显了其面向 LTL 规约的实用性。
\end{itemize}

本文结构如下：第 2 节回顾相关工作。第 3 节介绍预备知识。第 4 节介绍用于 LTL 验证的证明证书并与其他方法进行比较。第 5 节给出合成方法。第 6 节通过案例研究展示方法有效性。第 7 节总结全文。


\section{相关工作}\label{sec:related}
通过证明证书验证时序逻辑性质已经取得了显著进展。为了依据 LTL 规约验证离散时间系统，Wongpiromsarn 等~\cite{wongpiromsarn2015automata} 开发了一种状态三元组方法，该方法使用屏障证书阻断自动机中从初始状态到接受状态的所有简单路径，并要求枚举这些路径上的状态三元组 $(q_m, q_m', q_m'')$。Murali 等~\cite{murali2024closure} 引入了用于持久性验证的闭包证书，该证书为动力系统建立析取良基转移不变式 \cite{podelski2004transition}，并进一步将这一概念扩展到用于 $\omega$-正则规约的控制闭包证书~\cite{murali2025control}。此外，大量研究处理了特定时序行为的验证，例如可达性、安全性和避让~\cite{chakarov2013probabilistic,chatterjee2016algorithmic,chatterjee2024sound,chatterjee2024equivalence}。一些工作采用神经网络作为模板~\cite{wang2025verifying, nadali2024neural, zhao2020synthesizing} 来合成用于验证时序性质的证书，包括排名函数与归纳不变式的同时表示~\cite{giacobbe2025let}，以及针对 LTL 规约验证具有记忆的神经智能体~\cite{hosseini2025ltl}。


基于抽象的方法~\cite{vasile2016control, leahy2019control, kress2009temporal, coogan2015traffic, fainekos2009temporal, ding2014optimal, wongpiromsarn2009receding} 通过计算有限抽象，例如有限转移系统和马尔可夫决策过程（MDP），将 CPS 模型从连续域提升到离散域。我们注意到，基于抽象的技术已经用于针对 LTL 性质对连续空间动力系统进行验证和合成，例如~\cite{henzinger1997hytech, khaled2021omegathreads, lahijanian2011temporal, rungger2016scots} 中的结果。这些结果依赖于构建有限状态抽象，然后使用模型检测和合成技术~\cite{baier2008principles, tabuada2009verification}。总体而言，基于抽象的方法计算需求较高，并遭受维数灾难。无抽象方法，如证明证书，可以在一定程度上缓解计算负担。


%在我们的讨论背景下，我们假设已经构造了 LTL 规约 $\phi$ 的 NBA。从 LTL 公式到 NBA 的转换相对于公式大小具有指数复杂度~\cite{vardi2005automata}。构造 NBA 补自动机的复杂度已知为 EXPTIME~\cite{complementnba}。与标记广义 Büchi 自动机（LGBA）相比，将 LTL 规约转换为基于转移的广义 Büchi 自动机（TGBA）通常会得到更小的自动机~\cite{couvreur1999fly}。去广义化过程~\cite{giannakopoulou2002states} 可以将具有多个接受条件的 TGBA 转换为具有单一接受条件的基于转移的 Büchi 自动机（TBA）。存在多种工具可从 LTL 公式生成不同类型的自动机，例如 Owl~\cite{owl} 和 SPOT~\cite{spot}。
